\section{地域のエネルギー管理と評価の境界}
\label{theory:sustainability:energyterms}
再生可能エネルギーは太陽光、風力、水力、地熱、バイオマス等の資源を利用する。再生可能という性質と、設備の製造・輸送まで含めた環境影響ゼロは同じではない。CEMSは地域内の需要と供給を把握・調整する管理の枠組みであり、一つの建物だけの制御へ狭めない。

コージェネレーションは発電に伴う熱も有効利用する熱電併給である。熱需要がある場所・時刻との整合が効率を左右する\parencite{eth54_cogen2018}。発電効率だけと、熱を含む総合効率を混同しない。

ZEHは省エネと創エネにより年間の一次エネルギー収支を正味ゼロ以下とする住宅を扱う\parencite{eth54_zeh2018}。評価対象の用途、断熱・省エネの要件、計算の境界がある。毎時の購入電力がゼロであることや、家電を含むすべての使用エネルギーが常にゼロであることへ拡大しない。
